\documentclass[aspectratio=169,11pt]{beamer}

\usepackage[T1]{fontenc}
\usepackage[utf8]{inputenc}
\usepackage{lmodern}
\usepackage{amsmath,amssymb}
\usepackage{booktabs}
\usepackage{xcolor}
\usepackage{hyperref}
\usepackage{microtype}

\usetheme{Madrid}
\usecolortheme{default}

\definecolor{LaurierPurple}{HTML}{4B116F}
\definecolor{DeepBlue}{HTML}{0047AB}

\setbeamertemplate{navigation symbols}{}
\setbeamercolor{title}{fg=white,bg=LaurierPurple}
\setbeamercolor{frametitle}{fg=white,bg=LaurierPurple}
\setbeamercolor{block title}{fg=white,bg=DeepBlue}
\setbeamercolor{structure}{fg=DeepBlue}

\hypersetup{colorlinks=true,urlcolor=DeepBlue,linkcolor=DeepBlue}

\title[Regularization]{Module 3: Linear Models and Regularization}
\subtitle{EC 410K / EC 610I --- Machine Learning in Economics}
\author{M. Jahangir Alam, PhD}
\institute{Wilfrid Laurier University}
\date{Spring 2026}

\begin{document}

\begin{frame}[plain]
  \titlepage
\end{frame}

\begin{frame}{Module overview}
  \begin{itemize}
    \item High-dimensional regressors appear in modern empirical work: many controls, many baseline characteristics, many engineered features.
    \item \textbf{Ridge (L2)} stabilizes estimates via shrinkage; \textbf{Lasso (L1)} encourages sparsity (automatic variable selection).
    \item These tools are central to \textbf{prediction} and also appear inside modern \textbf{causal ML} pipelines as ways to estimate flexible nuisances (preview).
  \end{itemize}
\end{frame}

\begin{frame}{Learning objectives}
  \begin{enumerate}
    \item Write Ridge and Lasso penalized least squares objectives and interpret \(\lambda\).
    \item Explain why we standardize predictors before L1/L2 in practice.
    \item Tune \(\lambda\) with cross-validation and compare holdout error.
    \item Connect regularization to course papers: CNS (2022) and Athey--Wager (2021), plus replication examples with many covariates.
  \end{enumerate}
\end{frame}

\begin{frame}{Baseline problem: OLS in high dimensions}
  Minimize \(\|y - X\beta\|_2^2\). If \(p\) is large relative to \(n\), OLS can:
  \begin{itemize}
    \item overfit badly,
    \item be unstable (multicollinearity),
    \item produce models that generalize poorly.
  \end{itemize}
\end{frame}

\begin{frame}{Ridge regression (L2 penalty)}
  \[
    \widehat{\beta}^{\text{ridge}}
    = \arg\min_{\beta} \;\|y - X\beta\|_2^2 + \lambda \|\beta\|_2^2
  \]
  \begin{itemize}
    \item Shrinks coefficients toward 0; often improves out-of-sample prediction.
    \item Typically does not set coefficients exactly to 0.
  \end{itemize}
\end{frame}

\begin{frame}{Lasso (L1 penalty)}
  \[
    \widehat{\beta}^{\text{lasso}}
    = \arg\min_{\beta} \;\|y - X\beta\|_2^2 + \lambda \|\beta\|_1
  \]
  \begin{itemize}
    \item Can produce \textbf{exact zeros}: a form of variable selection.
    \item Tradeoff: useful for sparse signals; can be unstable with correlated predictors (elastic net is a common fix).
  \end{itemize}
\end{frame}

\begin{frame}{Choosing \(\lambda\)}
  \begin{itemize}
    \item \(\lambda\) controls bias--variance: larger \(\lambda\) \(\Rightarrow\) more shrinkage.
    \item Standard practice: grid search + \(K\)-fold cross-validation using out-of-sample loss.
  \end{itemize}
\end{frame}

\begin{frame}{Chernozhukov, Newey, and Singh (2022)}
  \begin{itemize}
    \item Develops \textbf{automatic debiased} machine learning approaches for valid inference on causal/structural parameters in high-dimensional settings.
    \item Big picture for students: first-stage ML must be handled carefully if you want confidence intervals for a low-dimensional target parameter.
  \end{itemize}
\end{frame}

\begin{frame}{Athey and Wager (2021): policy learning}
  \begin{itemize}
    \item Learns treatment assignment rules from observational data using robust objective functions.
    \item Connects to regularized/ML decision making: the goal is a \textbf{policy}, not only a coefficient.
  \end{itemize}
\end{frame}

\begin{frame}{Replication paper connections}
  \begin{itemize}
    \item \textbf{Anderson (2017):} many institutional covariates; ML regularization can help predict outcomes or estimate propensity models (later modules).
    \item \textbf{Goldin et al.\ (2016):} rich plan/market features; prediction of enrollment patterns vs causal claims about strategic pricing.
  \end{itemize}
\end{frame}

\begin{frame}{Python / Colab demonstration}
  \begin{enumerate}
    \item Simulate sparse \(\beta\) in \(p=60\) predictors.
    \item Fit Ridge and Lasso with \texttt{GridSearchCV}.
    \item Report holdout RMSE/R\(^2\) and visualize Lasso vs true \(\beta\).
  \end{enumerate}
\end{frame}

\begin{frame}[fragile]{Colab setup}
  \begin{verbatim}
!pip -q install numpy pandas scikit-learn matplotlib
  \end{verbatim}
\end{frame}

\begin{frame}{Student / group assignment}
  \begin{itemize}
    \item Explain Ridge vs Lasso with one geometric/intuitive analogy.
    \item Present CNS (2022) \textbf{or} Athey--Wager (2021) at overview level (students choose with instructor approval).
    \item Run the Colab script; discuss what would change if predictors were time-series autocorrelated.
    \item Propose a regularization-based extension for one replication paper (clear predictive vs causal scope).
  \end{itemize}
\end{frame}

\begin{frame}{Summary}
  \begin{itemize}
    \item Regularization makes linear models work better in high dimensions for \textbf{prediction}.
    \item Lasso induces sparsity; Ridge stabilizes; choose penalties with CV.
    \item Advanced causal inference with ML requires additional tools (orthogonalization, debiasing)---Module 4.
  \end{itemize}
\end{frame}

\begin{frame}[plain]
  \begin{center}
    \Large Questions?\\[0.8em]
    \normalsize Next: machine learning for causal inference (Module 4).
  \end{center}
\end{frame}

\end{document}
